\chapter{Der Zwischenspeicher ohne Ablaufzeit}\label{fall:cache}

\begin{merksatz}
\boxtitel{Die Vorlage}
\textbf{Zusammenhang.} Auswertung, Python. Referenzbereiche werden aus der Datenbank geholt und im Prozess gehalten.

\begin{center}
\begin{minipage}{0.9\textwidth}
\code{@lru\_cache(maxsize=None)}\\
\code{def bereich(groesse: GroessenCode, alter: int,}\\
\code{\ \ \ \ \ \ \ \ \ \ \ geschlecht: Geschlecht) -> Referenz:}\\
\code{\ \ \ \ return db.hole\_bereich(groesse, alter, geschlecht)}
\end{minipage}
\end{center}

\begin{center}
\begin{minipage}{0.88\textwidth}
\kursiv{Vier Zeilen, eine Anmerkung. Der Aufruf war vorher der teuerste Posten der Verlaufsansicht; seit dem Zwischenspeicher ist er nicht mehr messbar. Faktor vierhundert für ein Wort.}
\end{minipage}
\end{center}
\end{merksatz}

\begin{merksatz}
\boxtitel{Die Diagnose}
\begin{center}
{\small
\renewcommand{\arraystretch}{1.45}
\begin{tabular}{@{}P{3.0cm}P{2.2cm}P{2.6cm}@{}}
\toprule
\textbf{Befund} & \textbf{Schwere} & \\
\midrule
Keine Ablaufzeit & \textcolor{akzent}{\textbf{hoch}} & \textbf{bis Neustart}\\
Keine Größenbegrenzung & \textcolor{mittel}{mittel} & $\approx9000$ Einträge\\
\textbf{Nicht kohärent} & \textcolor{akzent}{\textbf{hoch}} & \textcolor{akzent}{\textbf{$6$ Prozesse}}\\
Nicht abschaltbar & \textcolor{mittel}{mittel} & \textcolor{akzent}{Fehlersuche}\\
\bottomrule
\end{tabular}}
\end{center}

\begin{center}
\begin{minipage}{0.88\textwidth}
\kursiv{Der erste Befund ist der klinisch bedeutsame: Ändert das Labor einen Referenzbereich, bewertet die Auswertung weiter nach dem alten --- bis jemand den Dienst neu startet. Das kann Wochen dauern.}
\end{minipage}
\end{center}

\textbf{Was tatsächlich geschah.}

\begin{center}
\renewcommand{\arraystretch}{1.45}
\begin{tabular}{@{}P{3.0cm}P{2.2cm}P{2.6cm}@{}}
\toprule
\textbf{Ereignis} & \textbf{Zeitpunkt} & \\
\midrule
Labor senkt obere Kaliumgrenze & \textbf{Mo, 8:00} & $5{,}1 \to 5{,}0$\\
Zwei Prozesse neu gestartet & Mo, 14:00 & \textcolor{akzent}{\textbf{übernehmen}}\\
Vier Prozesse alt & --- & \textcolor{akzent}{\textbf{bewerten anders}}\\
\textbf{Auffällig geworden} & \textcolor{akzent}{\textbf{Do}} & \textcolor{akzent}{\textbf{Nachfrage Station}}\\
\bottomrule
\end{tabular}
\end{center}

\begin{center}
\kursiv{Drei Tage lang bewerteten sechs Prozesse denselben Wert verschieden --- je nachdem, welcher die Anfrage bearbeitete. Ein Kaliumwert von $5{,}05$ war auf einem Bildschirm auffällig und auf dem nächsten nicht.}
\end{center}
\end{merksatz}

\begin{technik}
\boxtitel{Das Urteil: überarbeitet}
\begin{center}
\begin{minipage}{0.9\textwidth}
\code{\_bereiche = TTLCache(maxsize=20\_000, ttl=300)}\\
\code{\_sperre\ \ \ = threading.Lock()}\\
\\
\code{def bereich(g: GroessenCode, alter: int,}\\
\code{\ \ \ \ \ \ \ \ \ \ \ gs: Geschlecht) -> Referenz:}\\
\code{\ \ \ \ if not EINSTELLUNG.zwischenspeicher:}\\
\code{\ \ \ \ \ \ \ \ return db.hole\_bereich(g, alter, gs)}\\
\code{\ \ \ \ with \_sperre:}\\
\code{\ \ \ \ \ \ \ \ ...}
\end{minipage}
\end{center}

\begin{center}
\begin{minipage}{0.86\textwidth}
\kursiv{Drei Ergänzungen: eine Ablaufzeit von fünf Minuten, eine Größengrenze und ein Schalter. Der Zwischenspeicher bleibt im Prozess --- die Kohärenzfrage wird durch die Frist gelöst, nicht durch einen gemeinsamen Speicher.}
\end{minipage}
\end{center}
\end{technik}

\begin{merksatz}
\boxtitel{Die Begründung}
\begin{center}
{\small
\renewcommand{\arraystretch}{1.45}
\begin{tabular}{@{}P{2.8cm}P{2.0cm}P{2.2cm}P{2.0cm}@{}}
\toprule
& \textbf{Vorlage} & \textbf{Überarbeitet} & \\
\midrule
Zeilen & $4$ & $\mathbf{16}$ & \\
Trefferquote & $99{,}9\,\%$ & $\mathbf{99{,}5\,\%}$ & \\
Datenbanklast & $+0$ & \textcolor{mittel}{$+0{,}4\,\%$} & \\
\textbf{Veralterung} & \textcolor{akzent}{\textbf{unbegrenzt}} & \textcolor{gruen}{$\mathbf{\le 5}$ \textbf{min}} & \\
Prozesse uneins & \textcolor{akzent}{\textbf{tagelang}} & \textcolor{gruen}{$\le 5$ min} & \\
Abschaltbar & \textcolor{akzent}{nein} & \textcolor{gruen}{ja} & \\
Aufwand & --- & $1$ PT & \\
\bottomrule
\end{tabular}}
\end{center}

\begin{center}
\begin{minipage}{0.88\textwidth}
\kursiv{Vier Zehntel Prozent mehr Datenbanklast --- praktisch nichts --- gegen eine Veralterung von unbegrenzt auf fünf Minuten. Das ist der billigste Handel im ganzen Band.}
\end{minipage}
\end{center}

\textbf{Warum fünf Minuten?}

\begin{center}
\renewcommand{\arraystretch}{1.45}
\begin{tabular}{@{}P{2.4cm}P{2.2cm}P{2.2cm}P{2.0cm}@{}}
\toprule
\textbf{Frist} & \textbf{Treffer} & \textbf{Last} & \textbf{Veraltet}\\
\midrule
$30$ s & $94\,\%$ & \textcolor{akzent}{$+4{,}8\,\%$} & \textcolor{gruen}{$30$ s}\\
$\mathbf{5}$ \textbf{min} & $\mathbf{99{,}5\,\%}$ & \textcolor{gruen}{$+0{,}4\,\%$} & \textcolor{gruen}{$5$ min}\\
$1$ h & $99{,}9\,\%$ & \textcolor{gruen}{$+0{,}03\,\%$} & \textcolor{akzent}{$1$ h}\\
\bottomrule
\end{tabular}
\end{center}

\begin{center}
\kursiv{Die Kurve knickt bei etwa fünf Minuten: Darunter steigt die Last stark, darüber gewinnt man kaum noch. Die Zahl ist gemessen, nicht gewählt --- und das ist die eigentliche Arbeit an diesem Fall.}
\end{center}
\end{merksatz}

\begin{warnung}
\boxtitelr{Der Fehlschluss}
\begin{center}
\renewcommand{\arraystretch}{1.5}
\begin{tabular}{@{}P{3.2cm}P{2.2cm}P{2.4cm}@{}}
\toprule
\textbf{Naheliegend} & \textbf{Richtig?} & \\
\midrule
\enquote{Prozesseigene Speicher sind falsch} & \textcolor{akzent}{\textbf{nein}} & \textbf{mit Frist gut}\\
\enquote{Gemeinsamer Speicher wäre besser} & \textcolor{akzent}{\textbf{nein}} & \textcolor{akzent}{\textbf{$2000\times$ langsamer}}\\
\enquote{Bei Änderung ungültig machen} & \textcolor{mittel}{möglich} & \textcolor{akzent}{Meldung kann fehlen}\\
\textbf{\enquote{Jeder Zwischenspeicher braucht eine Frist}} & \textcolor{gruen}{\textbf{ja}} & \\
\bottomrule
\end{tabular}
\end{center}

\begin{center}
\begin{minipage}{0.88\textwidth}
\kursiv{Die zweite Zeile ist die teure Übertreibung: Ein gemeinsamer Speicher löst die Kohärenz und kostet zweihundert Mikrosekunden statt einer Zehntelmikrosekunde. Für einen Wert, der sich alle paar Monate ändert, ist die Frist die richtige Antwort.}
\end{minipage}
\end{center}

\begin{center}
\kursiv{Die dritte Zeile ist die technisch reizvollere und praktisch schwächere Lösung: Sie hängt daran, dass jede Änderungsmeldung ankommt --- die Annahme aus Band 23, Kapitel~2, die über eine Naht nicht zutrifft. Eine verlorene Meldung erzeugt einen dauerhaft falschen Speicher.}
\end{center}
\end{warnung}
\vspace{-4pt}
\noindent Klepp\-mann\zit{18} ordnet Ko\-hä\-renz und Ver\-al\-te\-rung ein; Gregg\zit{13} liefert das Vor\-ge\-hen, mit dem man die Frist misst statt sie zu wählen.
